\chapter{Vectores en el Servidor de Correo y Webmail}

En la infraestructura de red de las instituciones académicas a principios de los años 2000, los servidores de correo electrónico no solo procesaban el intercambio de mensajes mediante los protocolos SMTP, POP3 e IMAP, sino que integraban interfaces de acceso web (\textit{webmail}) desarrolladas comúnmente como aplicaciones CGI escritas en Perl o PHP. 

Estas aplicaciones web primitivas, ejecutadas típicamente bajo el usuario del servidor web o mediante envolturas \textit{suidperl}, adolecían con frecuencia de deficiencias estructurales: falta de cifrado HTTPS obligatorio, gestión insegura de identificadores de sesión en \textit{cookies} o parámetros URL, y vulnerabilidades de inyección de comandos.

En este capítulo se analiza la investigación pericial conducida sobre el servidor central de correo departamental (\texttt{victima3} / \texttt{io}), un examen que aportó la prueba más contundente sobre la intervención de un actor con conocimiento privilegiado del diseño interno de la red.

\section{Arquitectura de Seguridad del Servidor de Correo}

A diferencia de los servidores de cómputo \texttt{victima1} y \texttt{victima2}, el servidor de correo \texttt{victima3.foo.bar.mx} (identificado internamente en los registros como \texttt{io}) operaba bajo una política de bastionado (\textit{hardening}) deliberadamente más restrictiva:
\begin{itemize}
    \item Prestaba exclusivamente los servicios de red esenciales: recepción y entrega SMTP (puerto 25), descarga POP3 (puerto 110), consulta IMAP (puerto 143) e interfaz web \textit{OpenWebMail} sobre HTTP (puerto 80).
    \item **Denegación de acceso interactivo directo:** Ningún usuario ordinario tenía autorización ni acceso a un intérprete de comandos (\textit{shell}) remoto por SSH en esta máquina.
    \item **Esquema de administración restringido:** Para ejecutar tareas de administración, era mandatorio ingresar remotamente mediante una cuenta de salto altamente restringida (\texttt{admin}) y, una vez dentro del sistema, ejecutar el comando \texttt{su} suministrando la contraseña maestra de \texttt{root}.
\end{itemize}

Esta arquitectura de seguridad compartimentada era confidencial: en toda la universidad, la existencia de la cuenta de salto y el procedimiento específico de dos pasos para acceder a \texttt{root} era conocida exclusivamente por el administrador principal y por el titular de la cuenta \texttt{duende}.

\section{El Intento de Escalación en io: Evidencia de la Amenaza Interna}

Al auditar los registros de autenticación (\texttt{/var/log/authlog}) del servidor de correo correspondientes a la madrugada del 22 de abril, el equipo pericial localizó la secuencia crítica expuesta en el Listado~\ref{lst:io_authlog}.

\begin{lstlisting}[style=terminal,caption={Registros de autenticación en el servidor de correo io.},label={lst:io_authlog}]
Apr 18 22:37:42 io imapd[29834]: Login failed user=duende auth=duende host=dsl-198-51-100-131.isp-acceso.net [198.51.100.131]
Apr 18 22:37:45 io imapd[29834]: Login failed user=duende auth=duende host=dsl-198-51-100-131.isp-acceso.net [198.51.100.131]
Apr 22 06:18:01 io sshd[13560]: Accepted password for yyyyy from 192.0.2.20 port 33455 ssh2
Apr 22 06:18:08 io su: BAD SU yyyyy to root on /dev/ttyp0
Apr 22 06:18:18 io su: BAD SU yyyyy to root on /dev/ttyp0
Apr 22 06:18:47 io sshd[5397]: Failed password for duende from 192.0.2.20 port 33456 ssh2
Apr 22 06:18:52 io last message repeated 2 times
Apr 22 06:18:52 io sshd[5397]: Connection closed by 192.0.2.20
\end{lstlisting}

\subsection{Interpretación Forense de la Secuencia}

La reconstrucción técnica de este evento aportó indicios concluyentes:
\begin{enumerate}
    \item A las 06:18:01~h de la madrugada del 22 de abril, se registró una conexión SSH exitosa hacia la cuenta de usuario \texttt{yyyyy}, originada desde el servidor paritario \texttt{192.0.2.20} (\texttt{victima2}). El acceso fue limpio, sin intentos previos fallidos, lo que demuestra la posesión ilegítima de la clave correcta.
    \item Apenas siete segundos después (a las 06:18:08~h) y nuevamente a las 06:18:18~h, el operador intentó escalar privilegios a superusuario mediante el comando \texttt{su}, fallando en ambas ocasiones (\texttt{BAD SU yyyyy to root}). Este fallo confirma que el intruso no conocía la nueva contraseña de \texttt{root} instaurada preventivamente por el administrador días antes.
    \item Concurrente e inmediatamente después (a las 06:18:47~h), desde la misma máquina de origen (\texttt{victima2}), el atacante intentó iniciar una segunda sesión SSH buscando autenticarse como \texttt{duende}, cerrando la conexión tras registrar fallos sucesivos.
\end{enumerate}

La especificidad de este ataque descarta de plano a un intruso foráneo aleatorio: nadie fuera de la institución tenía conocimiento de la configuración interna de \texttt{io}, ni de los métodos de acceso privilegiado, ni de la relación de confianza entre \texttt{victima2} y el servidor de correo.

\section{Auditoría de Actividad en OpenWebMail}

La segunda línea de investigación sobre el servidor de correo abarcó las bitácoras de la aplicación de correo web \textit{OpenWebMail}. Esta interfaz registraba detalladamente las direcciones IP de los clientes, las sesiones autenticadas, los identificadores de mensajes leídos y las operaciones de buzón (traslados, supresión y reenvío).

El Listado~\ref{lst:openwebmail_curated} resume los accesos y operaciones ejecutadas durante el período bajo sospecha.

\begin{lstlisting}[style=terminal,caption={Extracto de operaciones recuperadas de la bitácora de OpenWebMail.},label={lst:openwebmail_curated}]
Fri Apr 11 12:25:26 2004 - [1219] (192.0.2.124) duende - movemsg - move 5 msgs from INBOX to saved-messages
Fri Apr 11 12:28:10 2004 - [1245] (192.0.2.124) duende - readmsg - id=<200404111728.i3BHSA612@victima1.com>
Tue Apr 15 07:09:53 2004 - [28749] (198.51.100.156) duende - logout - duende-session-0.12938
Fri Apr 18 19:35:10 2004 - [13101] (198.51.100.131) duende - readmsg - id=<4082B4F1.7060104@foo.bar.mx>
Fri Apr 18 19:38:25 2004 - [13129] (198.51.100.131) duende - logout - duende-session-0.88412
Mon Apr 21 21:10:04 2004 - [19401] (198.51.100.61)  duende - mail-trash - ids=<2403774.1049816@mail-agent>
\end{lstlisting}

\subsection{Análisis de las Operaciones en el Correo}

La auditoría de las bitácoras web evidenció que:
\begin{itemize}
    \item Las sesiones de \textit{OpenWebMail} del usuario \texttt{duende} se originaban alternadamente desde equipos de la red local universitaria (\texttt{192.0.2.124}) y desde el mismo conjunto de direcciones IP de acceso DSL residencial (\texttt{198.51.100.156}, \texttt{198.51.100.131}, \texttt{198.51.100.61}) identificadas en los ataques contra \texttt{victima1} y \texttt{victima2}.
    \item El usuario manipuló sistemáticamente sus buzones, trasladando mensajes confidenciales hacia carpetas secundarias (\texttt{saved-messages}) y purgando selectivamente registros de la papelera (\texttt{mail-trash}) para eliminar intercambios de mensajes ocurridos durante las fechas clave del compromiso.
    \item Se identificaron lecturas de mensajes administrativos institucionales donde se discutían anomalías de red en el departamento, lo que permitió al operador de la cuenta conocer que el personal técnico sospechaba de irregularidades en el centro de cómputo.
\end{itemize}

\section{Conclusiones del Análisis del Servidor de Correo}

La investigación en el servidor \texttt{victima3} / \texttt{io} cerró el círculo probatorio en torno a la doble naturaleza de la intrusión:
\begin{itemize}
    \item La correspondencia de direcciones IP externas entre las sesiones webmail de \texttt{duende} y los accesos de control a los bots IRC confirmó la unicidad del operador remoto.
    \item El intento fallido de escalamiento a superusuario en la madrugada del 22 de abril evidenció el móvil del actor interno: tras el cambio preventivo de claves ejecutado por los administradores, el atacante intentó desesperadamente reconquistar el control del servidor maestro mediante la explotación de cuentas secundarias y el conocimiento de la arquitectura interna de la red.
\end{itemize}
